\documentclass[oneside,doctype=proposal,fontset=windows,times=true]{tongjithesis}

% ============================================================
%  同 main.tex，两处区别：fontset=windows（宋体 SimSun / 黑体 SimHei）、times=true
%  （英文与数字用 Times New Roman，与官方 Word 一致；不开这项会回退到 TeX Gyre Termes，
%   字形相近但不是同一个字体，实测可查出差别）。
%
%  用途：在本机（已装微软中文字体）编出**与官方 Word 模板逐项对齐**的定稿 PDF。
%  ※ 这个文件在 Overleaf 上会编译失败（那边没有宋体/黑体），不要在那里选它当主文件。
%  编译：本机命令行 latexmk -xelatex main-windows.tex
% ============================================================

\begin{document}
\pagestyle{empty}

% ============================================================
%  文档类选项（在 main.tex 的 \documentclass[...]{tongjithesis} 中设置）
%  field=science    理工科（工科类、理科类，默认）— 章节编号：1 / 1.1 / 1.1.1
%  field=humanities 文科  （人文、法学、外语、艺术）        — 章节编号：一、/（一）/ 1.
%  经管类属社科：编号要求依理工，撰写依理工模板；26 届为过渡期，可选 science 或
%  humanities；自 27 届起统一使用 science。
% ============================================================

% ============================================================
%  封面信息
% ============================================================
% 以下字段全部取自已定稿的《2027届本科毕业论文任务书》（thesis/01_任务书/），
% 不另编造；如与最新任务书不一致，以任务书原件为准。
\school{电子与信息工程学院}
\major{人工智能}
\student{2352396}{禹尧珅}
\thesistitle{面向未来混合自动驾驶环境的风险传播机理研究}{}
\thesistitleeng{Research on Risk Propagation Mechanisms in Future Mixed Autonomous Driving Environments}{}
\thesisadvisor{李王明卉}
\thesisdate{\the\year}{\the\month}{\the\day}

\taskbookperiod{2026}{9}{14}{2027}{5}{26}

% 中期报告“专业审核意见”的同意/不同意复选框（仅编译 midterm.tex 时使用）：
% 默认两格均留空，供打印后手写勾选；如需提前在电子版中标注，取消以下一行注释
% \midtermdecision{agree} % 或 disagree

% ============================================================
%  信息说明页
% ============================================================

% 成果类型（三选一）：
%   thesis      — 毕业论文（摘要页标题：摘要）
%   design      — 毕业设计（软件/创意/建筑等非工程设计类；摘要页标题：摘要）
%   engineering — 毕业设计（工程设计类；摘要页标题：设计总说明）
\infotype{thesis}

% 内容简述（300字以内）
\infoabstract{%
请在此处填写论文内容简述，字数建议不超过300字。内容简述应概括论文的研究背景与问题动机、研究目标、主要方法与技术路线、实验设计与数据来源、关键实验结果及性能指标，以及研究结论与创新点，使读者能够快速了解本论文的核心贡献与研究意义。建议按照以下结构来组织简述内容：首先介绍研究背景与问题动机，充分说明研究的必要性与现有相关工作的主要不足之处；其次阐述所提出的方法或技术方案及其设计原理与主要实现细节；然后给出在典型数据集或实际应用场景下的主要实验结果与量化评价指标的全面横向对比分析；最后总结本研究的主要结论与当前局限性，并展望未来改进方向与潜在应用前景。请将以上示例文字全部替换为您的实际内容简述。
}

% 毕业设计：图纸数量与正文字数（成果类型为 design 或 engineering 时填写，两者须同时填写）
% \infodrawings{5}
% \infowordcount{15000}

% 毕业论文：正文总字数（成果类型为 thesis 时填写）
\infothesiswords{12345}

% 随附资料（每条用 \item 开头；不填则显示空白行）
\infomaterials{
  \item 随附资料名称一（如：全套图纸、程序源代码、计算书等）；
  \item 随附资料名称二
}

% ============================================================
%  中英文摘要页
% ============================================================
% 摘要页论文标题（可选）：设置摘要页展示的论文标题（不是“摘要”二字的标题）。未设置或两个参数都留空时，沿用 \thesistitle 与 \thesistitleeng。如需让摘要页使用不同标题或独立换行位置，取消以下两行注释。
% \abstracttitle{\TongjiThesis{}同济大学本科毕业设计（论文）模板与\LaTeX{}排版技术应用}{模板排版功能示例、学术写作规范与格式设置指南}
% \abstracttitleeng{\TongjiThesis{} template and \LaTeX{} typesetting for Tongji University undergraduate thesis}{A guide to typesetting features, academic writing standards, and formatting}

\MakeDocument
% ============================================================
%  开题报告正文骨架（AAD 毕设）
%
%  本文件当前是**框架**，不是成稿：
%    · 四节结构与两级标题完全照官方 2027 届模板，不动；
%    · 每个小节下用 \tjformslot 标出「这一段要写什么」，灰色、【待写】开头；
%    · 技术路线图用 \tjformfigslot 留图位，还没画；
%    · 参考文献只留位置，真实条目等检索完再填。
%
%  🔴 成稿前必做两件清理：
%     grep -n 'tjformslot\|tjformfigslot' chapters/proposal_body.tex   → 应为 0
%     grep -n 'tjformnote' chapters/proposal_body.tex                  → 应为 0（模板说明灰字，成稿删除）
%
%  🔴 防泄露（导师 2026-11 微信要求）：本文件不得写入 AAD 论文未公开的核心创新点——
%     消息可达性判据的具体定义、H1–H4 因果识别框架、十四类消息生命周期分类、
%     区域风险聚合的具体公式、以及现有量化结论。只写研究问题、技术路线、公开指标。
% ============================================================

\tjformsection{一、\tjheadertext 课题背景}
\tjformnote{根据实际情况写为毕业设计或毕业论文，下同（此项成稿删除）}

\tjformsubsection{1．选题背景与研究意义}
\tjformslot{开篇交代混合交通的必然性：自动驾驶、智能网联与人工智能技术发展，未来道路交通将长期是自动驾驶车辆与人工驾驶车辆共存。}
\tjformslot{再讲问题从哪来：两类车在感知能力、决策方式、响应时间、驾驶行为上差异明显，局部风险不再局限于单车或单一冲突，可能通过跟驰、制动、换道与交互行为向其他车辆和局部区域传播，并在特定条件下放大成级联响应。}
\tjformslot{然后指出缺口：只做单车风险独立分析，刻画不了整体安全演化；所以要研究局部风险如何产生、传播、累积，以及交通状态、车辆类型、交互关系对传播的影响。}
\tjformslot{最后落到意义：理解混合交通的安全机理、提高系统韧性，并为安全分析与运行管理提供理论与实验依据。}

\tjformsubsection{2．国内外研究现状（文献综述）}
\tjformslot{本节按四条线索组织，每条讲清「做了什么 — 局限在哪 — 本研究因此补什么」。真实文献检索完成后逐条填充，不得凭印象引用。}
\tjformslot{线索一 混合交通风险评价：替代安全指标（TTC、DRAC 等）的提出、适用条件与局限。}
\tjformslot{线索二 交通扰动与风险传播：扰动如何在车流中传播、现有研究做到什么尺度、缺什么。}
\tjformslot{线索三 微观交通仿真：SUMO 等平台在混合交通场景构建与校核中的用法与不足。}
\tjformslot{线索四 面向风险抑制的控制：现有干预思路与评价维度（安全性、效率、平稳性）。}
\tjformslot{小结：指出现有研究在「风险跨车辆、跨区域传播的可测与可干预」上的空白，由此引出本课题。}

\tjformsection{二、\tjheadertext 方案介绍}

\tjformsubsection{1．主要研究内容}
\tjformslot{与任务书「二、毕业论文应完成的主要内容」五条一一对应，每条一两句写清做什么、产出什么。}
\tjformslot{内容一 研究综述：调研混合交通、交通扰动传播、车辆交互风险、替代安全指标与交通仿真，梳理方法及适用范围，明确研究对象与技术路线。}
\tjformslot{内容二 风险测量方法：基于真实轨迹整理速度、加速度、车间距离与交互关系，选适用于典型跟驰冲突场景的安全评价指标，研究风险事件识别方法，给出指标合理取值范围。}
\tjformslot{内容三 仿真环境构建：基于 SUMO 等平台构建含自动驾驶与人工驾驶车辆的典型场景，设置车辆组成、交通流状态、驾驶行为与局部扰动，并用真实轨迹校核关键运动参数与风险状态。}
\tjformslot{内容四 传播规律与影响因素：设计不同交通状态、车辆组成与局部扰动条件下的实验，记录风险时空演化，分析传播、累积、衰减特征与影响范围、速度、强度。}
\tjformslot{内容五 抑制方法探索：依据传播规律研究利用车辆状态与风险信息提前干预的控制思路，对比安全性、效率与运行平稳性，评价抑制效果。}

\tjformsubsection{2．研究方案}
\tjformslot{本小节按「数据 — 方法 — 平台 — 实验 — 控制 — 路线图 — 可行性」七块展开。}
\tjformslot{数据：真实车辆轨迹数据（来源、字段、覆盖场景、预处理与异常缺失处理口径）。}
\tjformslot{方法：风险指标计算与风险事件识别；指标取值范围标定；不写论文未公开的判据定义。}
\tjformslot{平台：SUMO 版本、路网与场景构建方式、车辆行为参数设置、与真实轨迹的校核方法。}
\tjformslot{实验设计：交通状态、车辆组成、局部扰动三类因子如何组合；重复性与随机种子口径；评价指标与统计方法。}
\tjformslot{控制方法：干预信号从哪来、如何作用到车辆、对照设置与评价维度。}
\tjformfigslot{技术路线图（数据 → 指标 → 仿真 → 实验 → 控制，用代码画，图内不写核心算法细节）}
\tjformslot{可行性分析：数据已有、平台已有、方法有前期基础、时间安排可行（对应任务书进度）。}

\tjformsubsection{3．工作进度安排}
\tjformslot{照任务书「四、毕业论文的进度安排」七段填写，可整理成表；时间区间必须与任务书一致，不改。}

\tjformsection{三、\tjheadertext 主要参考文献}
\tjformrefnote{正文首行缩进2字符。参考文献悬挂0.74厘米，参考文献格式请参考GB/T 7714-2025《信息与文献 参考文献著录规则》，以下为样例。成稿请删除。}
\begin{tjformrefs}
\tjformslot{真实文献检索完成后填入，按 GB/T 7714-2025 著录，预计 15–20 条；每条必须能点开验证（DOI / arXiv 编号 / 期刊卷期页）。}
\end{tjformrefs}

\tjformsection{四、审核意见}

\begin{tjreviewbox}
  \tjformlabel{指导教师意见：}\tjformhint{（课题难度是否适中、工作量是否饱满、进度安排是否合理、工作条件是否具备、是否同意开题等）}

  \vspace{8\baselineskip}
  \tjformsignature{指导教师签名}{8em}
\end{tjreviewbox}

\begin{tjreviewbox}
  \tjformlabel{专业审核意见：}\tjformhint{（适当对各项内容给与意见和评价，字数不限。）}

  \vspace{6\baselineskip}
  \tjformsignature{负责人签名}{8em}
\end{tjreviewbox}


\end{document}
